(1) 经典单粒子配分函数 \(z_1=h^{-3}\int e^{-\beta p^2/(2m)}d^3p\int e^{-\beta\kappa r^2/2}d^3r=(2\pi mk_BT/h^2)^{3/2}(2\pi k_BT/\kappa)^{3/2}\)。\(Z_N=z_1^N/N!\)，仅在稀薄非简并条件下使用。
(2) \(n(r)=N(\kappa/(2\pi k_BT))^{3/2}e^{-\kappa r^2/(2k_BT)}\)，\(\langle r^2\rangle=3k_BT/\kappa\)。因 \(z_1\propto T^3\)，\(U=-\partial_\beta\ln Z_N=3Nk_BT\)，故 \(C_\kappa=3Nk_B\)。
(3) Stirling近似下 \(\ln Z_N=N\ln(z_1/N)+N\)，于是 \(S=k_B(\ln Z_N+\beta U)=Nk_B[\ln(z_1/N)+4]\)。绝热准静态过程保持 \(S,N\) 不变，因此 \(T^3\kappa^{-3/2}\) 不变，即 \(T/\sqrt\kappa\) 恒定；加紧谐阱会升温。
(4) 动量三个二次项贡献 \(3Nk_BT/2\)，坐标三个二次项亦贡献 \(3Nk_BT/2\)，合计与上式一致。\(\kappa\) 单位为能量每长度平方；经典近似要求 \(k_BT\gg\hbar\sqrt{\kappa/m}\) 且单粒子相空间占据远小于一。
